\SGASubsection{1.7}{The pinned groups \(Z_{\Delta_1}\)}
\SGARefTarget{sga3:XXIII:1.7}

Let \(\Delta_1\) be a subset of the system \(\Delta\) of simple roots
of the pinned group \(G\). Let \(T_{\Delta_1}\) be the maximal torus
of \(\bigcap_{\alpha\in\Delta_1}\ker(\alpha)\), and set
\[
  Z_{\Delta_1}=\operatorname{Centr}_G(T_{\Delta_1}).
\]
Put \(R'=\mathbb Z\Delta_1\cap R\). It is known
(\SGARefLink{sga3:xxii:para:5.10.7}{Exposé~XXII, 5.10.7}) that \(Z_{\Delta_1}\) is a reductive \(S\)-group
with radical \(T_{\Delta_1}\), that \((T,M,R')\) is a splitting of
this group, and that \(\Delta_1\) is a system of simple roots. Thus
\[
  (Z_{\Delta_1},T,M,R',\Delta_1,(X_\alpha)_{\alpha\in\Delta_1})
\]
is a pinned \(S\)-group. We shall always give \(Z_{\Delta_1}\) this
pinning. In particular, we shall consider
\[
  Z_\alpha=Z_{\{\alpha\}},\qquad
  Z_{\alpha\beta}=Z_{\{\alpha,\beta\}}.
\]
Write \(B_{\Delta_1}=B\cap Z_{\Delta_1}\). It is known
(\textit{loc.\ cit.}) that this is the canonical Borel subgroup
\NDE{That is, the Borel subgroup of \(Z_{\Delta_1}\) corresponding to
\(R'^+=R'\cap R^+\).}
of \(Z_{\Delta_1}\), and that its unipotent part is
\(U_{\Delta_1}=U\cap Z_{\Delta_1}\). In particular,
\[
  B_\alpha=T\cdot U_\alpha.
\]
We write
\[
  U_{\alpha\beta}=U_{\{\alpha,\beta\}}
  =U\cap Z_{\alpha\beta}
  =\prod_{\gamma\in R_{\alpha\beta}^{+}}U_\gamma,
\]
where \(R_{\alpha\beta}^{+}\) is the set of positive roots that are
linear combinations of \(\alpha\) and \(\beta\).

Now let \(f:G\to G'\) be a morphism of pinned \(S\)-groups. If
\(d:R\to R'\) is the corresponding bijection and
\(\Delta_1\subset\Delta\), then
\(d(\Delta_1)=\Delta'_1\subset\Delta'\). Clearly \(f\) maps
\(T_{\Delta_1}\) into \(T'_{\Delta'_1}\), and hence
\(Z_{\Delta_1}\) into \(Z'_{\Delta'_1}\). The corresponding
\(S\)-morphism
\[
  f_{\Delta_1}:Z_{\Delta_1}\longrightarrow Z'_{\Delta'_1}
\]
is compatible with the canonical pinnings. It defines a morphism of
pinned root data
\[
  \mathscr R(f_{\Delta_1}):
  \mathscr R(Z_{\Delta_1},T,M,\ldots)
  \longrightarrow
  \mathscr R(Z'_{\Delta'_1},T',M',\ldots),
\]
and the underlying morphisms \(M'\to M\) of \(\mathscr R(f)\) and
\(\mathscr R(f_{\Delta_1})\) coincide.

\SGASubsection{1.8}{Study of the group
\(\operatorname{Norm}_G(T)\)}
\SGARefTarget{sga3:XXIII:1.8}

For every pair \((\alpha,\beta)\) of simple roots, let
\(n_{\alpha\beta}\) be the order of \(s_\alpha s_\beta\) in the Weyl
group \(W\). In particular \(n_{\alpha\alpha}=1\), and therefore
\((w_\alpha w_\beta)^{n_{\alpha\beta}}\in T(S)\).

\SGARefTarget{sga3:XXIII:1.8.1}
\SGAStatement{Definition 1.8.1}{For every
\((\alpha,\beta)\in\Delta\times\Delta\), set
\[
  t_{\alpha\beta}=(w_\alpha w_\beta)^{n_{\alpha\beta}}\in T(S).
\]
Moreover, set (\SGARefLink{sga3:xx:3.1}{Exposé~XX, 3.1})
\[
  t_\alpha=t_{\alpha\alpha}=w_\alpha^2=\alpha^\ast(-1)\in T(S).
\]}

\SGARefTarget{sga3:XXIII:1.8.2}
\SGAStatement{Proposition 1.8.2}{Let \(H\) be an \(S\)-functor in
groups that carries direct sums of schemes to products (for example,
a sheaf for the Zariski topology). Let
\[
  f_T:T\longrightarrow H
\]
be a morphism of groups, and let \(h_\alpha\in H(S)\)
(\(\alpha\in\Delta\)). There exists a necessarily unique morphism of
groups
\[
  f_N:\operatorname{Norm}_G(T)\longrightarrow H
\]
inducing \(f_T\) on \(T\) and satisfying \(f_N(w_\alpha)=h_\alpha\)
for \(\alpha\in\Delta\) if and only if the following two conditions
hold:
\begin{enumerate}
  \item[\SGAInlineRefTarget{sga3:XXIII:1.8.2:item:i}(i)] for every \(\alpha\in\Delta\),
    \[
      f_T(s_\alpha(t))=h_\alpha f_T(t)h_\alpha^{-1}
    \]
    for every \(t\in T(S')\), \(S'\to S\); equivalently,
    \(f_T\circ s_\alpha=\operatorname{int}(h_\alpha)\circ f_T\);
  \item[\SGAInlineRefTarget{sga3:XXIII:1.8.2:item:ii}(ii)] for every
    \((\alpha,\beta)\in\Delta\times\Delta\),
    \[
      f_T(t_{\alpha\beta})
        =(h_\alpha h_\beta)^{n_{\alpha\beta}}.
    \]
\end{enumerate}}

\SGAProof
Equip \((\mathrm{Sch})\) with the topology \(\mathcal T\) generated
by the pretopology whose covering families are direct sums. The
hypothesis says that \(H\) is a \(\mathcal T\)-sheaf. Let \(L\) be
the free group on the generators \((m_\alpha)_{\alpha\in R}\), and
let \(L_1\) be the invariant subgroup generated by
\[
  (m_\alpha m_\beta)^{n_{\alpha\beta}},
  \qquad(\alpha,\beta)\in\Delta\times\Delta.
\]
Let \(g:L\to W\) be defined by \(g(m_\alpha)=s_\alpha\). It is known
(\SGARefLink{sga3:XXI:5.1}{Exposé~XXI, 5.1}) that \(g\) induces an isomorphism
\(\overline g:L/L_1\xrightarrow{\sim}W\). Let \(L\) act on \(T\)
through \(g\), equivalently by \(m_\alpha\cdot t=s_\alpha(t)\).
Writing \(L_S\) for the constant group defined by \(L\), form the
semidirect product \(N=T\cdot L_S\). There is a unique morphism of
\(S\)-groups
\[
  h:N=T\cdot L_S\longrightarrow\operatorname{Norm}_G(T)
\]
such that \(h(m_\alpha)=w_\alpha\) and \(h(t)=t\) for every
\(t\in T(S')\), \(S'\to S\).

Let \(N_1\) be the invariant subgroup sheaf of \(N\) generated by
\[
  t_{\alpha\beta}^{-1}(m_\alpha m_\beta)^{n_{\alpha\beta}},
  \qquad(\alpha,\beta)\in\Delta\times\Delta.
\]
Clearly \(N_1\subset\ker h\), so \(h\) induces
\[
  h_1:N/N_1\longrightarrow\operatorname{Norm}_G(T).
\]
We prove that \(h_1\) is an isomorphism. Since \(h|_T\) is the
canonical immersion, hence a monomorphism, the canonical morphism
\(T\to N/N_1\) is also a monomorphism and induces an isomorphism
\(T\xrightarrow{\sim}TN_1/N_1\). Thus \(h_1\) induces an
isomorphism \(TN_1/N_1\xrightarrow{\sim}T\). It remains only to show
that
\[
  h_2:N/TN_1\longrightarrow\operatorname{Norm}_G(T)/T
\]
is an isomorphism. But \(TN_1\) is the invariant
\(\mathcal T\)-subgroup sheaf of \(N\) generated by \(T\) and the
\((m_\alpha m_\beta)^{n_{\alpha\beta}}\); hence it is the subgroup
sheaf generated by \(T\) and \(L_1\), namely \(T\cdot(L_1)_S\).
Therefore \(h_2\) identifies with
\[
  \overline g:L_S/(L_1)_S\longrightarrow W_S,
\]
which is an isomorphism by construction.

The proof of the proposition is now immediate. Necessity is clear.
Conversely, condition \SGARefLink{sga3:XXIII:1.8.2:item:i}{(i)} gives a morphism
\[
  u:N\longrightarrow H
\]
such that \(u(m_\alpha)=h_\alpha\) for \(\alpha\in\Delta\) and
\(u|_T=f_T\). Condition \SGARefLink{sga3:XXIII:1.8.2:item:ii}{(ii)} says that \(u\) vanishes on \(N_1\),
and the result follows.

\SGASubsection{1.9}{Faithfulness of the functor \(\mathscr R\)}
\SGARefTarget{sga3:XXIII:1.9}

\SGARefTarget{sga3:XXIII:1.9.1}
\SGAStatement{Proposition 1.9.1}{The functor \(\mathscr R\) of \SGARefLink{sga3:XXIII:1.6}{1.6} is
faithful: if
\[
  f,g:G\rightrightarrows G'
\]
are two morphisms of pinned groups such that
\(\mathscr R(f)=\mathscr R(g)\), then \(f=g\).}

Indeed, \(f\) and \(g\) agree on \(T\), on every \(U_\alpha\)
(\(\alpha\in\Delta\)), and on every \(U_{-\alpha}\)
(\(\alpha\in\Delta\)); it therefore suffices to apply:

\SGARefTarget{sga3:XXIII:1.9.2}
\SGAStatement{Lemma 1.9.2}{Let \(S\) be a scheme, let \(G\) be a
pinned reductive \(S\)-group, and let \(H\) be an \(S\)-presheaf in
groups separated for the fppf topology. Let
\[
  f,g:G\rightrightarrows H
\]
be two morphisms of \(S\)-groups. The following conditions are
equivalent:
\begin{enumerate}
  \item[\SGAInlineRefTarget{sga3:XXIII:1.9.2:item:i}(i)] \(f=g\);
  \item[\SGAInlineRefTarget{sga3:XXIII:1.9.2:item:ii}(ii)] \(f\) and \(g\) agree on \(T\), on every \(U_\alpha\)
    (\(\alpha\in\Delta\)), and on every \(U_{-\alpha}\)
    (\(\alpha\in\Delta\));
  \item[\SGAInlineRefTarget{sga3:XXIII:1.9.2:item:iii}(iii)] \(f\) and \(g\) agree on \(T\) and on every
    \(U_\alpha\) (\(\alpha\in\Delta\)), and
    \(f(w_\alpha)=g(w_\alpha)\) for every \(\alpha\in\Delta\).
\end{enumerate}}

\SGAProof
The implication \SGARefLink{sga3:XXIII:1.9.2:item:i}{(i)}\(\Rightarrow\)\SGARefLink{sga3:XXIII:1.9.2:item:ii}{(ii)} is trivial, while
\SGARefLink{sga3:XXIII:1.9.2:item:ii}{(ii)}\(\Rightarrow\)\SGARefLink{sga3:XXIII:1.9.2:item:iii}{(iii)} follows immediately from the definition of
\(w_\alpha\) in \SGARefLink{sga3:XXIII:1.2}{1.2}. It remains to prove
\SGARefLink{sga3:XXIII:1.9.2:item:iii}{(iii)}\(\Rightarrow\)\SGARefLink{sga3:XXIII:1.9.2:item:i}{(i)}. If \(\alpha\in R\), there is a sequence
\(\{\alpha_i\}\subset\Delta\) such that
\[
  \alpha=s_{\alpha_1}s_{\alpha_2}\cdots
    s_{\alpha_n}(\alpha_{n+1}),
\]
and hence
\[
  U_\alpha
    =\operatorname{int}(w_{\alpha_1}\cdots w_{\alpha_n})
       U_{\alpha_{n+1}} .
\]
Thus \(f\) and \(g\) agree on every \(U_\alpha\). They consequently
agree on \(\Omega\), and therefore agree everywhere
(\SGARefLink{sga3:xxii:prop:4.1.11}{Exposé~XXII, 4.1.11}).

\SGARefTarget{sga3:XXIII:1.9.3}
\SGAStatement{Remark 1.9.3}{If \(G\) is semisimple, the hypothesis
that \(f\) and \(g\) agree on \(T\) may be omitted in \SGARefLink{sga3:XXIII:1.9.2:item:ii}{(ii)} and \SGARefLink{sga3:XXIII:1.9.2:item:iii}{(iii)}.
Indeed, \(G\), as an fppf sheaf, is generated by the \(U_\alpha\),
\(\alpha\in R\) (\SGARefLink{sga3:xxii:rem:6.2.2:item:a}{Exposé~XXII, 6.2.2(a)}).}
